\subsection{Experimentos com \textit{Hybrid Pipelines}}

Os experimentos com \textit{hybrid pipelines}\footnote{Disponível em: \url{https://github.com/giovannabbottino/hybrid-pipelines}.} foram realizados por meio de um serviço que combina um modelo de linguagem de grande porte (\textit{Large Language Model} -- LLM) com evidências provenientes da Wikidata\footnote{Disponível em: \url{https://www.wikidata.org/wiki/Wikidata:Main_Page}.}. Nessa abordagem, as entidades e os conceitos identificados no texto podem ser associados a recursos de uma base de conhecimento estruturada, cujas informações são utilizadas durante a construção do grafo de conhecimento.

Diferentemente do experimento baseado em \textit{few-shot prompting}, o fluxo híbrido incorpora informações provenientes de uma fonte externa durante o processo de construção do grafo. As menções identificadas no texto são associadas, quando possível, a recursos da Wikidata, e as informações recuperadas são utilizadas como evidências para a geração da representação em RDF/Turtle. Além dessas evidências, o fluxo utiliza procedimentos locais de processamento, validação e correção da representação produzida.

\begin{figure}[htbp]
    \centering
    \caption{Fluxo de criação de grafos de conhecimento com \textit{Hybrid Pipelines}.}
    \includegraphics[width=0.65\linewidth]{metodo/figuras/experimento-c.jpg}
    \label{fig:experimento-c}
    \source{da autora.}
\end{figure}

Conforme apresentado na Figura~\ref{fig:experimento-c}, o fluxo experimental é organizado em quatro etapas principais: \textit{(i)} extração das menções presentes no texto; \textit{(ii)} enriquecimento das entidades com informações provenientes da Wikidata; \textit{(iii)} validação e, quando necessário, correção da representação em RDF/Turtle; e \textit{(iv)} salvamento do grafo de conhecimento resultante.

A construção ocorre de forma sequencial. Inicialmente, as entidades e os conceitos são extraídos do texto. Em seguida, essas menções são associadas, quando possível, a recursos da Wikidata e enriquecidas com informações recuperadas da base. Os dados obtidos são então utilizados na construção da representação em RDF/Turtle, que passa por procedimentos de validação e correção antes de ser encaminhada à etapa final de salvamento.


\subsubsection{Definição dos \textit{prompts}}

Nesse experimento, foram utilizados três \textit{prompts} com funções complementares: um \textit{prompt} de sistema, um \textit{prompt} destinado à extração das menções e um \textit{prompt} destinado à construção da representação em RDF/Turtle.

O \textit{prompt} de sistema, apresentado no Apêndice~\ref{apendiceC}, estabelece o comportamento geral do modelo nas chamadas realizadas pelo fluxo. Entre suas instruções, determina-se que o LLM produza exclusivamente o formato solicitado pelo \textit{prompt} específico da tarefa, evitando a inclusão de introduções, explicações, exemplos, notas adicionais ou delimitadores de \textit{markdown}.

O \textit{prompt} de extração, apresentado no Apêndice~\ref{apendiceD}, recebe o texto de entrada e solicita ao modelo a identificação das entidades e dos conceitos relevantes. A resposta é estruturada em formato JSON e utilizada como entrada para as etapas posteriores de resolução e enriquecimento das entidades.

O \textit{prompt} de construção, apresentado no Apêndice~\ref{apendiceE}, recebe um \textit{payload} formado pelo texto original, pelas entidades identificadas e pelas relações recuperadas durante o enriquecimento com a Wikidata. Com base nessas informações, o modelo é orientado a construir a representação do grafo de conhecimento em RDF/Turtle, preservando os identificadores recuperados na etapa anterior.


\subsubsection{Extração de menções}

A primeira etapa consiste na identificação das entidades e dos conceitos presentes no texto de entrada. Para essa finalidade, o serviço utiliza o modelo de linguagem em conjunto com o \textit{prompt} específico de extração.

Para cada menção identificada, são previstos os atributos \texttt{surface}, correspondente à forma textual encontrada no documento; \texttt{start}, referente à posição inicial da ocorrência; \texttt{end}, correspondente à posição final; \texttt{entity\_type}, que representa o tipo semântico atribuído à menção; e \texttt{confidence}, utilizado para registrar o valor de confiança associado à identificação.

De forma simplificada, a estrutura solicitada ao modelo pode ser representada por:

\begin{verbatim}
{
  "entities": [
    {
      "surface": "...",
      "start": 0,
      "end": 5,
      "entity_type": "Entity",
      "confidence": 0.0
    }
  ]
}
\end{verbatim}

Após o recebimento da resposta do LLM, o conteúdo é interpretado estritamente como JSON e as menções identificadas são realinhadas às respectivas ocorrências no texto original. Esse procedimento permite associar cada entidade ou conceito ao trecho efetivamente encontrado no documento. Uma resposta malformada, vazia ou sem menções utilizáveis interrompe a requisição.

Depois de uma extração válida, o conjunto de menções pode ser complementado por regras aplicadas ao próprio texto para padrões lexicais específicos. Essas regras integram a etapa híbrida, mas não substituem uma extração inválida ou vazia.

Em seguida, as menções são deduplicadas com base na forma textual e nas respectivas posições inicial e final. Dessa forma, ocorrências distintas de uma mesma expressão podem ser preservadas quando aparecem em diferentes posições do texto ou representam sentidos distintos.

Por fim, aplica-se o limite máximo de menções definido para o experimento. O resultado dessa etapa corresponde ao conjunto de entidades e conceitos encaminhado à etapa de enriquecimento com a Wikidata.


\subsubsection{Enriquecimento com a Wikidata}

A segunda etapa consiste na associação das menções extraídas a recursos disponíveis na Wikidata e na recuperação de evidências estruturadas relacionadas às entidades identificadas.

Cada menção obtida na etapa anterior é utilizada como termo de busca para recuperar possíveis candidatos exclusivamente pelo Wikidata MCP. Os recursos encontrados são posteriormente comparados para a seleção daquele que melhor representa a entidade mencionada no texto. Uma falha do MCP é propagada sem consulta a outra fonte.

Quando a resolução é realizada, a menção textual é associada a um identificador no formato \texttt{Q-ID}, fornecendo uma referência estruturada à entidade representada no grafo. Após a seleção do candidato, são recuperadas informações relacionadas ao respectivo recurso, incluindo declarações da Wikidata que podem ser utilizadas na construção da representação final.

O serviço também verifica a existência de relações diretas entre as entidades resolvidas no mesmo texto. Quando uma declaração associada a determinado recurso aponta para outra entidade também identificada no conjunto analisado, essa relação é preservada como evidência estruturada.

Ao final dessa etapa, são obtidos três componentes principais: as menções identificadas no texto, os recursos da Wikidata associados a essas menções e as relações diretas encontradas entre as entidades resolvidas.

Quando uma menção não pode ser associada a um identificador da Wikidata, ela pode permanecer no conjunto processado como uma entidade não resolvida. Assim, a ausência de um \texttt{Q-ID} não implica sua exclusão das etapas posteriores.

As informações produzidas nessa fase são utilizadas na construção do RDF/Turtle, permitindo que o LLM receba, além do texto original, as evidências estruturadas recuperadas da Wikidata.


\subsubsection{Construção e validação do RDF/Turtle}

A terceira etapa compreende a construção e a validação estrita da representação em RDF/Turtle.

Antes da chamada ao modelo, é criado um \textit{payload} JSON contendo o texto original, a atribuição da fonte, as entidades identificadas e as relações recuperadas durante o enriquecimento. De forma simplificada, o conteúdo apresenta a seguinte estrutura:

\begin{verbatim}
{
  "text": "...",
  "source_attribution": "Source: Wikidata",
  "entities": [...],
  "relationships": [...]
}
\end{verbatim}

Esse conteúdo substitui, em tempo de execução, o marcador \texttt{\$\{PAYLOAD\}} presente no \textit{prompt} de construção.

O \textit{prompt} orienta o modelo a preservar os identificadores das entidades previamente resolvidas. Assim, os recursos associados à Wikidata devem ser representados no formato \texttt{wd:Q...}, utilizando os identificadores recuperados na etapa de enriquecimento.

As entidades também devem ser acompanhadas por \texttt{rdfs:label}. Para os recursos resolvidos, prioriza-se como rótulo a forma superficial presente no texto analisado. As relações são materializadas como triplas dirigidas e os predicados são representados no \textit{namespace} \texttt{kg:}. Quando uma relação estruturada é recuperada durante o enriquecimento, essa informação pode ser utilizada na definição da conexão representada no grafo.

Nos casos em que o conteúdo textual apresenta uma negação, o \textit{prompt} prevê a utilização da propriedade:

\begin{verbatim}
kg:negated true
\end{verbatim}

Após a geração, a representação produzida pelo modelo passa por uma etapa de limpeza. São removidos eventuais delimitadores de blocos de código, conteúdos anteriores à primeira declaração \texttt{@prefix} e notas adicionais que não façam parte da sintaxe RDF/Turtle.

Em seguida, o conteúdo é submetido à biblioteca \texttt{rdflib} para análise sintática no formato Turtle. Para ser considerada válida, a representação deve ser interpretável pelo analisador e conter pelo menos uma tripla.

Quando a representação é considerada válida, o fluxo segue para a etapa de salvamento. Caso contrário, nenhum procedimento local altera ou preserva parcialmente a saída.

Caso a representação não seja válida e ainda existam tentativas disponíveis, uma nova chamada à mesma etapa do LLM pode ser realizada, utilizando o conteúdo inválido e as informações referentes ao erro encontrado durante a análise. O número de tentativas é limitado. Se nenhuma tentativa produzir uma representação válida, o processamento é interrompido sem mecanismo substituto.

Além das relações recuperadas diretamente da Wikidata, o pós-processamento pode acrescentar relações heurísticas entre entidades identificadas no mesmo segmento textual. Essas relações são produzidas localmente pela implementação e são, portanto, distintas das declarações recuperadas da Wikidata.


\subsubsection{Salvamento do grafo de conhecimento}

A quarta e última etapa consiste no salvamento do grafo de conhecimento após sua validação.

Conforme apresentado na Figura~\ref{fig:experimento-c}, somente um RDF validado diretamente no modo primário selecionado segue para essa etapa.

A representação RDF/Turtle validada constitui o resultado final do processo de construção. O grafo preserva as entidades identificadas no texto, os identificadores da Wikidata recuperados durante o enriquecimento, os respectivos rótulos e as relações materializadas durante a geração e o pós-processamento.

Além do RDF/Turtle, o serviço mantém informações produzidas nas etapas anteriores, como as entidades identificadas, as relações recuperadas e a saída produzida pelo LLM durante a extração. Esses elementos permitem relacionar o grafo final às evidências utilizadas em sua construção.

De forma simplificada, o resultado produzido pelo fluxo pode ser representado por:

\begin{verbatim}
{
  "text": "...",
  "entities": [...],
  "relationships": [...],
  "rdf": "...",
  "source_attribution": "Source: Wikidata",
  "llm": {
    "entity_extraction": "..."
  }
}
\end{verbatim}

Dessa forma, a quarta etapa consolida os resultados das fases anteriores e encerra o fluxo de construção do grafo de conhecimento.


\subsubsection{Configuração do experimento}

A Tabela~\ref{tab:llm-config-b} apresenta os principais parâmetros utilizados no experimento com \textit{hybrid pipelines}. Assim como no experimento baseado em \textit{few-shot prompting}, são apresentados o nome do parâmetro, seu tipo, sua função e o valor utilizado.

\begin{table}[htbp]
\centering
\caption{Configurações utilizadas no experimento com \textit{Hybrid Pipelines}.}
\label{tab:llm-config-b}
\small
\renewcommand{\arraystretch}{1.2}

\begin{tabular}{
p{3.2cm}
p{1.5cm}
p{6.0cm}
p{2.0cm}
}
\toprule
\textbf{Parâmetro} &
\textbf{Tipo} &
\textbf{Descrição} &
\textbf{Valor utilizado} \\
\midrule

\texttt{stream}
&
\textit{boolean}
&
Desabilita o envio incremental da resposta
&
\texttt{false}
\\

\texttt{seed}
&
\textit{integer}
&
Semente utilizada durante a geração
&
\texttt{42}
\\

\texttt{temperature}
&
\textit{float}
&
Controla o grau de aleatoriedade da geração
&
\texttt{0}
\\

\texttt{top\_k}
&
\textit{integer}
&
Limita a seleção aos \textit{tokens} mais prováveis
&
\texttt{40}
\\

\texttt{top\_p}
&
\textit{float}
&
Define o limiar da amostragem \textit{nucleus}
&
\texttt{0.9}
\\

\texttt{min\_p}
&
\textit{float}
&
Define a probabilidade mínima considerada
&
\texttt{0.05}
\\

\texttt{num\_predict}
&
\textit{integer}
&
Número máximo de \textit{tokens} produzidos pelo modelo
&
\texttt{1536}
\\

\texttt{timeout seconds}
&
\textit{integer}
&
Tempo máximo de espera da requisição, em segundos
&
\texttt{180}
\\

\texttt{candidate\_limit}
&
\textit{integer}
&
Número máximo de candidatos recuperados para cada menção
&
\texttt{3}
\\

\texttt{entity\_limit}
&
\textit{integer}
&
Número máximo de menções encaminhadas à resolução
&
\texttt{16}
\\

\bottomrule
\end{tabular}
\end{table}

A Tabela~\ref{tab:llm-config-b} apresenta apenas os parâmetros diretamente relacionados ao comportamento das principais etapas do experimento. Configurações relacionadas à infraestrutura, como endereços dos serviços, caminhos utilizados para armazenamento dos registros, tempos de espera das requisições e políticas de repetição, não foram incluídas por não constituírem parâmetros centrais da configuração experimental.

No perfil utilizado, até 16 menções são encaminhadas à etapa de resolução das entidades. As chamadas realizadas ao Ollama utilizam semente \texttt{42}, temperatura \texttt{0} e limite de \texttt{1536} \textit{tokens} de saída. Quando os parâmetros da requisição não são modificados, o LLM dispõe de três tentativas para gerar o RDF; uma saída inválida após esse limite encerra a requisição.


\subsubsection{Considerações sobre a reprodutibilidade}

A utilização de uma semente fixa e de temperatura igual a zero busca reduzir a variabilidade das respostas produzidas pelo modelo de linguagem. Entretanto, esses parâmetros, isoladamente, não constituem uma verificação de determinismo da geração.

Além das características inerentes ao modelo, o fluxo depende de informações recuperadas de uma base de conhecimento externa. Consequentemente, alterações posteriores na Wikidata podem modificar os candidatos, as declarações e as relações recuperadas em diferentes momentos de execução.

A construção do RDF é sempre realizada pelo LLM. Uma falha de validação não provoca reparo local nem troca para um mecanismo alternativo; somente as tentativas adicionais da mesma etapa do LLM são permitidas.

Por essa razão, os registros produzidos durante a execução são relevantes para a rastreabilidade do experimento, pois conservam informações sobre as chamadas realizadas ao modelo, as entidades identificadas, as relações recuperadas e o RDF/Turtle obtido ao final do processamento.

Em síntese, o experimento com \textit{hybrid pipelines} combina o modelo de linguagem, informações provenientes da Wikidata MCP e procedimentos locais de processamento e validação. O LLM é utilizado na extração das menções e na geração da representação RDF/Turtle; a Wikidata fornece candidatos, identificadores e relações estruturadas; e as rotinas locais realizam o processamento e a validação sem corrigir a saída inválida. O grafo somente é encaminhado à etapa final após a obtenção de uma representação RDF/Turtle considerada válida.
